\chapter{Tiến và Lùi}
\markboth{Tiến và Lùi}{Advance and Retreat}

\section{Một người lùi một bước để tránh làm tổn thương thêm.}
\begin{truyen}{Một người lùi một bước để tránh làm tổn thương thêm.}{Stepping back can be wise, not weak.}
\chuhoa{K}{hi cuộc cãi vã lên cao, anh biết nếu nói tiếp chỉ có thể khiến mọi thứ vỡ thêm. Anh im, rời khỏi cuộc nói chuyện một lúc, rồi quay lại khi cả hai đã bình tĩnh hơn. Lùi đúng lúc không phải thua; đó là cách giữ phần người trong lúc nóng nhất.}
\textit{(When the argument rose too high, he knew that continuing would only deepen the damage. He went silent, stepped away for a while, and returned when both sides were calmer. Stepping back at the right moment is not defeat; it is a way of protecting humanity in the hottest moment.)}
\end{truyen}
\begin{baihoc}
Không phải lùi nào cũng là hèn; có cái lùi là trí khôn.
\textit{(Not every retreat is cowardice; some retreats are wisdom.)}
\end{baihoc}
\begin{ghinhoanh}
\textbf{Short English to remember:}

A wise step back can save much.
\end{ghinhoanh}
\begin{tuhocnhanh}
1. Em có tình huống nào đáng lẽ nên dừng sớm hơn? \textit{(Is there a situation where you should have stopped earlier?)}

2. Với em, lùi lúc nào là khôn chứ không phải yếu? \textit{(For you, when is retreat wise rather than weak?)}
\end{tuhocnhanh}
\ngancach

\section{Một học sinh tiến chậm nhưng không lùi khỏi mục tiêu.}
\begin{truyen}{Một học sinh tiến chậm nhưng không lùi khỏi mục tiêu.}{Slow advance still counts as advance.}
\chuhoa{E}{m chưa thể giỏi ngay, nhưng mỗi tuần em hiểu thêm một mảng nhỏ. Người khác không thấy bước nhảy lớn nên tưởng em đứng yên. Chỉ em biết mình không lùi, và nhiều khi điều quan trọng nhất chính là đừng bỏ đường.}
\textit{(The student could not become excellent immediately, but each week one small area became clearer. Others saw no dramatic leap and assumed there was no progress. Only the student knew there had been no retreat, and sometimes the most important thing is simply not leaving the road.)}
\end{truyen}
\begin{baihoc}
Tiến thật không phải lúc nào cũng ồn ào.
\textit{(Real progress is not always noisy.)}
\end{baihoc}
\begin{ghinhoanh}
\textbf{Short English to remember:}

Quiet progress is still progress.
\end{ghinhoanh}
\begin{tuhocnhanh}
1. Em có đang xem nhẹ những bước nhỏ của mình không? \textit{(Are you underestimating your small steps?)}

2. Điều gì cho thấy em vẫn đang tiến dù chưa thấy lớn? \textit{(What shows that you are still moving forward even without big results?)}
\end{tuhocnhanh}
\ngancach

\section{Lùi khỏi một đám đông để tiến gần hơn đến chính mình.}
\begin{truyen}{Lùi khỏi một đám đông để tiến gần hơn đến chính mình.}{Leaving the crowd can restore your direction.}
\chuhoa{C}{ó thời gian em tham gia mọi cuộc vui chỉ vì sợ bị bỏ lại. Nhưng càng đi cùng đám đông, em càng thấy mình rỗng. Khi bớt chen vào những nơi không hợp, em lại có thời gian đọc, học và nghe rõ lòng mình hơn.}
\textit{(For a time the student joined every gathering out of fear of being left out. Yet the more the crowd was followed, the emptier it felt inside. Once those unsuitable spaces were left behind, time returned for reading, learning, and hearing the inner voice more clearly.)}
\end{truyen}
\begin{baihoc}
Có những cái lùi khỏi đám đông chính là cái tiến vào chiều sâu.
\textit{(Some retreats from the crowd are advances into depth.)}
\end{baihoc}
\begin{ghinhoanh}
\textbf{Short English to remember:}

Leave the crowd, find your path.
\end{ghinhoanh}
\begin{tuhocnhanh}
1. Đám đông nào đang kéo em ra khỏi điều quan trọng? \textit{(Which crowd is pulling you away from what matters?)}

2. Nếu lùi ra, em sẽ tiến được gần hơn với điều gì? \textit{(If you step back, what could you move closer to?)}
\end{tuhocnhanh}
\ngancach

\section{Cứ lao tới mà không biết lúc nào nên dừng.}
\begin{truyen}{Cứ lao tới mà không biết lúc nào nên dừng.}{Forward is not always upward.}
\chuhoa{A}{nh quen với suy nghĩ phải luôn đi tiếp, luôn nhận thêm, luôn thắng thêm. Vì thế anh không nhận ra khi sức khỏe đã xuống, gia đình đã xa, và niềm vui đã mòn. Có người tưởng mình đang tiến vì chưa từng lùi, nhưng thật ra đang lao ra khỏi điều đáng sống.}
\textit{(He was used to thinking he must always move on, take more, and win more. Because of that, he failed to see that health had weakened, family had grown distant, and joy had worn thin. Some people think they are advancing because they never retreat, yet they are rushing away from what makes life worth living.)}
\end{truyen}
\begin{baihoc}
Không phải cứ tiến là đúng hướng.
\textit{(Moving forward is not always moving rightly.)}
\end{baihoc}
\begin{ghinhoanh}
\textbf{Short English to remember:}

Forward is not always better.
\end{ghinhoanh}
\begin{tuhocnhanh}
1. Em đang tiến về đâu, và cái giá là gì? \textit{(Where are you moving toward, and what is the cost?)}

2. Em có cần dừng lại để kiểm tra hướng đi không? \textit{(Do you need to pause and check your direction?)}
\end{tuhocnhanh}
\ngancach

\section{Lùi để sửa mình rồi tiến chắc hơn.}
\begin{truyen}{Lùi để sửa mình rồi tiến chắc hơn.}{A pause for repair can strengthen the next step.}
\chuhoa{S}{au nhiều lần làm việc cẩu thả, cô xin nghỉ một thời gian ngắn để học lại kỹ năng và sắp xếp thói quen. Người khác tưởng cô tụt lại, nhưng khi trở lại, cô làm việc điềm tĩnh và bền hơn hẳn. Có những bước lùi chỉ là quãng lấy đà cho một lần tiến tử tế.}
\textit{(After many careless mistakes, she took a short break to relearn skills and reorder habits. Others thought she had fallen behind, but when she returned, her work was steadier and more durable. Some retreats are simply the space needed for a healthier advance.)}
\end{truyen}
\begin{baihoc}
Lùi để sửa không làm giảm giá trị; nó làm tăng chất lượng của bước kế tiếp.
\textit{(Stepping back to repair does not reduce value; it improves the quality of the next step.)}
\end{baihoc}
\begin{ghinhoanh}
\textbf{Short English to remember:}

Repair before you rush.
\end{ghinhoanh}
\begin{tuhocnhanh}
1. Em có phần nào cần tạm chậm hoặc lùi để sửa? \textit{(Is there any area where you need to slow down or step back in order to repair?)}

2. Bước tiến nào của em sẽ tốt hơn nếu được chuẩn bị lại từ gốc? \textit{(Which step forward would improve if prepared again from the roots?)}
\end{tuhocnhanh}
\ngancach
